% Generated by analyze_allbench_problem_boundary.py; do not hand-edit values.
\begin{table*}[t]
\centering
\footnotesize
\setlength{\tabcolsep}{3pt}
\caption{Five-workload result. P/H and $T_N/T_P$ are paired geomean timing ratios. Identical H/P/N placement digests and copy totals make a row a timing control: its ratio cannot be attributed to a locality-placement decision.}
\label{tab:allbench-problem-boundary}
\begin{tabularx}{\textwidth}{@{}lrr>{\raggedright\arraybackslash}p{0.13\textwidth}r>{\raggedright\arraybackslash}X@{}}
\toprule
Workload & P/H & P copy reduction & H/P/N placement & $T_N/T_P$ & Interpretation \\
\midrule
FDTD & 1.064x & 0.0\% & identical & 0.922x & scheduler-path/order/system timing control \\
CG & 1.005x & 0.0\% & identical & 1.001x & scheduler-path/order/system timing control \\
Cholesky & 1.405x & 85.8\% & different & 0.920x & opportunity; effective; conservative gate cost \\
LU & 1.173x & 38.8\% & different & 1.208x & opportunity; effective; over-optimization counterexample \\
MiniWeather & 0.982x & 0.0\% & identical & 0.991x & scheduler-path/order/system timing control \\
\bottomrule
\end{tabularx}
\vspace{2pt}\par\scriptsize Positive copy reduction is the independent modeled-opportunity signal. $T_N/T_P>1$ favors prefix safety; $<1$ favors exact gate removal. Copy quantities are scheduler-model outputs, not measured link traffic. Timing ratios in identical-placement rows quantify run-order/system variation only.
\end{table*}
